\chapter{Predictive Coding, Representation Comparison, and Adaptive Computation}
\label{ch:related-work}

This chapter serves two distinct functions. It first bounds the literature context and
then introduces the dissertation's own theoretical constructs. The separation prevents
external PC/BP results from being mistaken for evidence supporting the dissertation's
C01--C11 claims.

The literature is organized around the three final research questions. Backpropagation
is first treated as the baseline computational training line. The discussion then
moves from the neurobiological formulation of predictive coding (PC) to contemporary
inference-based learning variants and the question of how their updates relate to
backpropagation (BP). Methods for comparing internal representations and gradients are
considered separately. Finally, work on adaptive computation, early exit, and selective
prediction is used to position QWake as a compute-continuation decision problem under
a registered admissibility criterion.

Several claim types must remain distinct even when they appear together in the
literature: similarity of final quality, similarity of internal representations,
similarity of gradient directions, exact equality of updates under special conditions,
and computational efficiency. This dissertation keeps these levels separate and uses
each source only within the question it directly addresses.

\section{Backpropagation as the Baseline Computational Line}
\label{sec:related-bp}

The classical work of Rumelhart, Hinton, and Williams established practical training
of multilayer networks by backpropagation and linked changes in hidden representations
to reduction of output error \cite{rumelhart1986backprop}. In modern terms, BP can be
viewed as an efficient implementation of the chain rule for derivatives through a
composition of functions. In this dissertation, BP is not treated as a biological
model but as a well-defined reference algorithmic regime against which paired
comparisons can be constructed.

This positioning matters for interpretation. Matching or similar quality metrics do
not imply matching internal learning trajectories. Similarity of one layer's gradient
does not automatically transfer to other layers, and an approximation to BP under
special conditions does not specify the cost of satisfying those conditions. BP is
therefore used as a behavioral baseline, a source of reference gradients, and a point
of comparison for matched profiling.

\section{Hierarchical Predictive Coding}
\label{sec:related-pc}

The historical predictive-coding line includes models in which neighboring hierarchy
levels exchange predictions and prediction errors. Rao and Ballard connected this
architecture to visual processing and hierarchical explanation of sensory input
\cite{rao1999predictive}. Prediction-error units have an explicit computational role,
and network state includes internal variables that change during iterative inference,
not only parameters.

For machine learning, Whittington and Bogacz are particularly relevant
\cite{whittington2017approximation}. They compare supervised learning in a PC network
using local Hebbian updates with BP and show that, under specified conditions, PC
updates approximate backpropagation updates. The problem is thereby narrowed from a
general comparison of biological principles to an algorithmic question: under which
dynamics, initial conditions, and update rules does a local inference procedure
reproduce or approximate a reference gradient?

Contemporary work also changes and accelerates the internal inference dynamics. Hybrid
Predictive Coding combines multiple temporal inference regimes and motivates separation
of fast and slower computational processes \cite{tschantz2023hybrid}. This line is
important here because PC iterativity makes the number of internal computation steps a
first-class quantity and therefore permits separate study of intermediate-state
sufficiency and the price of further refinement.

\section{Relations Between Predictive Coding and Backpropagation}
\label{sec:related-pc-bp}

The PC/BP relation cannot be reduced to one universal equivalence claim. Different
works use different dynamics, initialization schemes, and update orders, so every
result must retain its applicability domain.

Whittington and Bogacz analyze approximation of BP by local rules
\cite{whittington2017approximation}. Song et al. identify conditions under which a
predictive-coding network reproduces BP updates exactly \cite{song2020exact}.
Salvatori et al. extend exact-BP reproduction to convolutional and recurrent networks,
which is directly relevant to the convolutional architecture used here
\cite{salvatori2021convolutional}. Millidge, Tschantz, and Buckley generalize the
analysis to arbitrary differentiable computation graphs and describe asymptotic
convergence of PC gradients toward BP gradients in their construction
\cite{millidge2022arbitrary}.

These results do not remove the need to state comparison constraints. Rosenbaum
analyzes several PC/BP regimes in which equivalence statements have different strengths
\cite{rosenbaum2022relationship}. The present work also incorporates the official 2025
correction \cite{rosenbaum2025correction}; the original article is not interpreted
without the corrected version.

A separate line concerns the computational cost of exact or BP-like PC variants.
Zahid, Guo, and Fountas study modern PC variants as a possible neuromorphic alternative
and derive temporal-complexity limitations for the investigated constructions
\cite{zahid2023critical}. This is especially relevant to RQ2: locality of updates or a
theoretical relation to BP does not itself establish a computational advantage.

The literature on PC efficiency and scaling has subsequently broadened. Salvatori et
al. show that changing the temporal order of weight updates in iPC improves stability
and efficiency relative to standard PC in their studied tasks
\cite{salvatori2024stable}. Pinchetti et al. make efficiency and scalability explicit
benchmarking targets, introduce PCX, and evaluate substantially larger architectures
and task sets than typical early PC experiments \cite{pinchetti2025benchmarking}.
Innocenti, Achour, and Buckley show that the specialized $\mu$PC parameterization can
stably train PC networks deeper than one hundred layers in their investigated
conditions while exposing scaling pathologies of the standard parameterization
\cite{innocenti2025mupc}. Goemaere et al. in ePC examine the mismatch between local PC
state dynamics and their digital simulation and propose an error-based parameterization
for faster and more scalable computation \cite{goemaere2026epc}.

These works neither replace the measurements of this dissertation nor explain its
results retrospectively. They show that computational cost, signal propagation with
depth, and scaling are independent open coordinates of PC. Mechanistic attribution
and matched cost profiling are therefore treated here as separate empirical tasks
after behavioral and gradient comparability have been assessed.

The literature thus supports the sequential organization of the present work: first
establish \emph{what} agrees or differs between regimes, then localize the region of
difference, test its registered mechanistic attribution, and only then discuss the
computational price of a concrete implementation.

\section{Similarity of Internal Representations and Gradients}
\label{sec:related-representations}

Comparing neural networks only by accuracy or loss leaves open whether their internal
representations are similar. SVCCA combines singular-value decomposition and canonical
correlation analysis to compare activations across networks or training times
\cite{raghu2017svcca}. Morcos, Raghu, and Bengio extend CCA-based representation
analysis and discuss robustness of similarity conclusions \cite{morcos2018insights}.

Kornblith et al. systematically compare representation-similarity measures and
introduce centered kernel alignment (CKA) as a useful metric
\cite{kornblith2019similarity}. CKA is the primary representation diagnostic here,
while CCA-based methods provide complementary literature context.

Representation similarity and gradient similarity answer different questions. The
former describes activation geometry on a selected sample; the latter describes the
local parameter-update direction. Stage~3A therefore does not aggregate these into a
single ``similarity score''. Gradient cosine similarity, gradient-norm ratio, and CKA
are complementary coordinates. This distinction is particularly important for regimes
that can preserve representation geometry while strongly changing early-layer gradient
magnitude.

\section{Theoretical Framework of the Study}
\label{sec:theory-framework}

The theoretical constructs are introduced in dependency order rather than experimental
execution order. Task-relative equivalence and PC-TREF come first; the distinct
mechanistic PC-CATM level follows; their relation is then fixed; and only afterward is
QWake-PC introduced as a compute-action selection architecture. None of these steps
turns PC-CATM into a component of PC-TREF, and PC-TREF sufficiency never becomes an
economic-benefit claim.

The literature on PC/BP relations primarily asks when two computational procedures
produce equal or similar gradients and updates. Compute-budget control requires more.
Two numerically different states may require the same computational action; conversely,
states close under a simple diagnostic may require different actions if further
computation materially changes the required final response.

The dissertation therefore uses a bounded three-level construction. PC-TREF formalizes
which state differences must be preserved for an admissible compute-control action.
PC-CATM models correction formation in predictive coding and thereby supplies
mechanistically interpretable sources of diagnostic features. QWake-PC links these
levels to action and budget selection. QWake-FP is only one preregistered bounded
instantiation for FixedPred.

The construction is not claimed as a universal theory of predictive coding. Action
spaces, norms, thresholds, relations, and sufficiency criteria are defined only
relative to the registered task, Torch2PC implementation, and experimental domain.

\subsection{Relation to State Abstraction and Action-sufficient Representations}
\label{sec:theory-state-abstraction-basis}

The general principle that a representation should preserve only state differences
capable of changing the decision is not new to this dissertation. MDP state-abstraction
theory studies partitions that preserve optimal-control properties
\cite{li2006stateabstraction}; bisimulation metrics provide a quantitative basis for
state similarity and aggregation while controlling changes in value
\cite{ferns2004metrics}. More directly, Huang et al. introduce action-sufficient
representations (ASRs), minimal representations retaining the information required for
subsequent action selection \cite{huang2022actionsufficient}. Raman, Milani, and Fang
formulate a related decision-preserving state abstraction in which a common conceptual
representation must not merge states with different optimal actions
\cite{raman2026decisionrelevant}.

PC-TREF does not claim novelty for this general principle and is not a new MDP,
bisimulation, or representation-learning theory. Its contribution here is specialization
of action-relative sufficiency to residual iterative PC computation. The finite action
set $\mathcal A$, required partition $q_R$, diagnostic partition $q_I$, domain
$\Omega$, and tolerance $\delta_R$ are registered in advance. Diagnostics may use only
pre-action information, while sufficiency is audited with a post-action reference.
Required-partition defect and common-admissibility defect are separated explicitly,
and representation suitability for compute control is separated both from the
admissibility of a concrete decision rule and from the cost of producing and using the
representation. PC-TREF is therefore a registered PC specialization of existing
state-abstraction and action-sufficiency ideas, not a claim to have invented
decision-preserving abstraction itself.

\subsection{PC-TREF: Equivalence Relative to the Computational Task}
\label{sec:theory-pc-tref}

PC-TREF is the Predictive-Coding Task-Relative Equivalence Framework. Its premise is
that compute control does not require equality of complete internal states; it requires
preserving state differences that can change the admissible action.

Let $x\in\mathcal X$ be an iterative-inference state and
$\phi_I(x)\in\mathcal Y_I$ a selected diagnostic representation. If a discrete
diagnostic partition is required, a frozen mapping $q_I$ defines
\[
  x\sim_I y
  \quad\Longleftrightarrow\quad
  q_I\!\left(\phi_I(x)\right)=q_I\!\left(\phi_I(y)\right).
\]
This relation induces a diagnostic quotient: within one class, the controller
intentionally does not distinguish states under the selected representation.

For continuous features, exact equality is often impractical. An operational
closeness relation can instead be defined as
\[
  x\approx_{I,\varepsilon}y
  \quad\Longleftrightarrow\quad
  d_I\!\left(\phi_I(x),\phi_I(y)\right)\leq\varepsilon_I,
\]
with metric, normalization, and tolerance frozen before analysis. Thresholded closeness
is not itself treated as an equivalence relation because it can be non-transitive and
does not induce a quotient without a separate partition rule.

The second PC-TREF object is the required equivalence, defined by the task rather than
the diagnostic. Let $\mathcal A$ be the preregistered action set and $L_R(a,x)$ the
loss associated with choosing action $a$ in state $x$. The dissertation uses
\[
  \operatorname{Regret}_R(a;x)
  =L_R(a,x)-\inf_{b\in\mathcal A}L_R(b,x)\geq0.
\]
It measures how much worse $a$ is than the best action within the registered set
$\mathcal A$. Action $a$ is admissible when the regret does not exceed the frozen
tolerance $\delta_R$.

The required partition $q_R$ is part of the registered task contract and encodes the
action distinctions that must be preserved. PC-TREF neither asserts uniqueness of
$q_R$ nor derives a minimal partition outside the registered domain from $L_R$,
$\mathcal A$, and $\delta_R$. Let $\mathcal C_R$ be the set of required classes and
$d_R:\mathcal C_R\rightarrow\mathcal A$ the registered action assignment. Define
\[
  \mathcal T_R=(\mathcal A,L_R,\delta_R,q_R,d_R).
\]
Then
\[
  x\sim_R y\quad\Longleftrightarrow\quad q_R(x)=q_R(y),
\]
and each class must satisfy
\[
  \sup_{x:q_R(x)=c}
  \operatorname{Regret}_R\!\left(d_R(c);x\right)\leq\delta_R.
\]
The existence of another action admissible for several classes does not permit their
merger. This matters when a universal fallback action exists: exact continuation may be
admissible almost everywhere, but it must not collapse all of $\Omega$ into one class
and thereby erase the question of early-action admissibility.

QWake-FP therefore specializes the required partition directly by early-action
admissibility:
\[
 q_R^{\mathrm{QW}}(x)=
 \begin{cases}
   \texttt{EARLY\_ADMISSIBLE},&
   \operatorname{Regret}_R(a_{\mathrm{early}};x)\leq\delta_R,\\
   \texttt{CONTINUE\_REQUIRED},&\text{otherwise}.
 \end{cases}
\]
For this specialization,
$d_R(\texttt{EARLY\_ADMISSIBLE})=a_{\mathrm{early}}$ and
$d_R(\texttt{CONTINUE\_REQUIRED})=a_{\mathrm{continue}}$; the contract separately
requires the assigned action to be admissible within each class. Canonical continuation
remains the exact fallback, but is not used to merge the two required classes. PC-TREF
therefore preserves the distinction between early-action admissibility and later
economic benefit.

The key PC-TREF comparison is between diagnostic and required partitions. Let
\[
  E_I^q=\{(x,y):q_I(\phi_I(x))=q_I(\phi_I(y))\},
  \qquad
  E_R^q=\{(x,y):q_R(x)=q_R(y)\}.
\]
The exact required-partition defect is
\[
  \mathfrak D_{I\to R}^{q}=E_I^q\setminus E_R^q.
\]
It contains pairs intentionally merged by the diagnostic representation although the
registered task assigns them to different required decision classes. Another common
fallback action does not remove this defect; universal \texttt{CONTINUE} therefore
does not permit merging \texttt{EARLY\_ADMISSIBLE} with
\texttt{CONTINUE\_REQUIRED}.

Separately define common admissibility
\[
  A_R^{\delta}
  =\left\{(x,y):\exists a\in\mathcal A:\;
  \operatorname{Regret}_R(a;x)\leq\delta_R
  \land
  \operatorname{Regret}_R(a;y)\leq\delta_R\right\}.
\]
This relation need not be transitive and is not used as a second equivalence relation.
Diagnostic pairs with no common admissible action form the narrower defect
\[
  \mathfrak D_{I\to R}^{q,\delta}=E_I^q\setminus A_R^{\delta}.
\]
Thus, $\mathfrak D_{I\to R}^{q}$ records loss of required discrimination, while
$\mathfrak D_{I\to R}^{q,\delta}$ records at least one diagnostic pair lacking a
common admissible action. The pairwise property does not replace a class-level test for
a single common admissible action across three or more states. The two defect notions
need not coincide.

For a registered domain $\Omega\subseteq\mathcal X$, strong sufficiency of a discrete
diagnostic representation is defined through the required partition:
\[
  \phi_I\text{ is sufficient relative to }(\mathcal T_R,\Omega)
  \quad\Longleftrightarrow\quad
  E_I^q|_{\Omega}\subseteq E_R^q|_{\Omega}
  \quad\Longleftrightarrow\quad
  \mathfrak D_{I\to R}^{q}|_{\Omega}=\varnothing.
\]
When the superscript is unnecessary, write $E_I:=E_I^q$ and $E_R:=E_R^q$. The
criterion is relative to the frozen $q_I$, $q_R$, and $\Omega$ and does not establish
sufficiency outside them. A partition defect is not yet a dangerous action of a
concrete rule. A conservative rule can respond to lost discrimination by continuing,
reducing coverage or adding computation; a dangerous accept occurs only when it
proposes the early action for a state requiring continuation.

Hence $\phi_I$ is sufficient relative to the registered required partition not because
it reconstructs the full internal state, but because it does not merge states from
different required classes on the registered domain. Admissibility of a concrete rule
is checked separately. Minimality is not claimed. Only a predefined representation
family is compared, and a more compact representation can be preferred only after
action admissibility and feature-computation cost are considered.

Full sufficiency $E_I|_\Omega\subseteq E_R|_\Omega$ is a strong property of the entire
diagnostic partition. A selective controller can test a narrower property without
claiming full sufficiency of $\phi_I$. For action $a$, define its admissibility set
\[
  S_a=\{x\in\Omega:\operatorname{Regret}_R(a;x)\leq\delta_R\}.
\]
In the remainder of the dissertation, the phrase ``a state task-sufficient for action
$a$'' means only $x\in S_a$ and is not synonymous with full sufficiency of the
diagnostic representation. For a rule $\pi$ proposing $a$ on
\[
  A_{\pi}(a)=\{x\in\Omega:\pi(x)=a\},
\]
call the rule selectively sufficient for action $a$ on $\Omega$ when
\[
  A_{\pi}(a)\subseteq S_a.
\]
The corresponding rule-admissibility defect is
\[
  \mathfrak D_{\pi,a}=A_{\pi}(a)\setminus S_a.
\]
This is one-sided: the rule must be correct where it proposes the specialized action,
but may abstain elsewhere even when the same action is admissible. Non-zero selective
coverage therefore does not prove full representation sufficiency, and partition-defect
absence remains distinct from rule admissibility.

In QWake-FP, $a_{\mathrm{early}}$ denotes the registered analytic-completion action.
The experimental dangerous-accept counter is defined later in C2 methodology. At the
theoretical level, zero observed dangerous accepts means no observed elements of
$\mathfrak D_{\pi,a_{\mathrm{early}}}$ on the finite calibration surface. It is only
an empirical manifestation of
$A_{\pi}(a_{\mathrm{early}})\subseteq S_{a_{\mathrm{early}}}$ and does not prove the
inclusion on all of $\Omega$, establish
$\mathfrak D_{I\to R}^{q}=\varnothing$ for the full representation, or replace
independent confirmation.

\subsection{PC-CATM: Mechanism of Correction Formation and Cancellation}
\label{sec:theory-pc-catm}

PC-TREF specifies \emph{which} state distinctions matter for a decision but does not
explain \emph{where} informative differences come from. PC-CATM, the
Predictive-Coding Correction Aggregation and Transport Model, provides that
mechanistic level.

For layer $l$, denote prediction, state, and prediction error by
\[
  p_l=f_l(h_{l-1};\theta_l),\qquad e_l=p_l-h_l.
\]
With a symmetric positive-semidefinite precision operator $\Pi_l\succeq0$, the local
quadratic energy is
\[
  F_l=\frac12 e_l^{\mathsf T}\Pi_l e_l,\qquad F=\sum_j F_j.
\]
For internal state $h_l$, the two neighboring energy terms give the structural
decomposition
\[
  -\frac{\partial F}{\partial h_l}
  =\Pi_l e_l-J_{h,l+1}^{*}\Pi_{l+1}e_{l+1}.
\]
PC-CATM does not claim this standard energy-derived dynamic as a new PC algorithm. Its
role is to treat structurally distinct terms as diagnostic channels and analyze their
activity, transport, and cancellation separately. The self channel is
\[
  c_l^{(\mathrm{self})}=\Pi_l e_l,
\]
and the transported upper channel is
\[
  c_l^{(\mathrm{upper})}
  =-J_{h,l+1}^{*}\Pi_{l+1}e_{l+1},
\]
where $J_{h,l+1}^{*}$ is the adjoint operator transporting the influence of upper-layer
error to the current state. The resulting correction is the channel aggregate
\[
  u_l=\mathcal S_l(\mathbf c_l)=\sum_k c_l^{(k)}.
\]

A small $\|u_l\|$ can arise through at least two mechanisms. The individual channels
may themselves be small, indicating a region of low activity, or individually large
channels may cancel. PC-CATM keeps these cases separate and therefore does not treat a
small resulting residual as automatic evidence that further computation is unnecessary.

Define total channel activity and resulting-correction norm
\[
  \alpha_l=\sum_k\|c_l^{(k)}\|,
  \qquad
  m_l=\left\|\sum_k c_l^{(k)}\right\|,
\]
and for $\alpha_l>0$ the residual-correction coefficient
\[
  \chi_l=\frac{m_l}{\alpha_l}.
\]
Small $\chi_l$ with non-zero $\alpha_l$ indicates strong mutual cancellation, but
remains diagnostic rather than a stopping permission. PC-CATM also considers
inter-layer transport: a non-zero upper-layer error component may weakly affect a lower
state because of Jacobian-transport geometry.

The normative glossary names exact and thresholded instances of these phenomena NCZ,
ECZ, and TNZ. For the dissertation, the important distinction is among low source
activity, cancellation of active channels, and suppression during transport. They
supply different possible PC-TREF features and must not be collapsed into one scalar
``state error''.

Every PC-CATM norm and threshold exists only inside a measurement contract specifying
space, norm, scale, numerical precision, device, layer, step, and aggregation rule. A
numerically small value in one layer or regime therefore does not automatically transfer
to another. PC-CATM is a mechanistic feature model and does not make a compute-control
decision by itself.

In glossary terms, channel norms and pairwise interactions form \emph{correction
geometry}. The exact correction-zero set is the kernel of the channel-summation
operator,
\[
  \mathcal Z_l^{(0)}=\ker\mathcal S_l.
\]
Its trivial part
\[
  \mathrm{NCZ}_l^{(0)}=\{\mathbf0\}
\]
corresponds to all canonical channels being exactly zero, while the non-trivial part
\[
  \mathrm{ECZ}_l^{(0)}=\ker\mathcal S_l\setminus\{\mathbf0\}
\]
corresponds to exact cancellation of non-zero channels. Exact NCZ and ECZ therefore
partition $\ker\mathcal S_l$ but do not include every state with a small non-zero
resulting correction.

The numerical protocol does not identify exact sets with threshold classes. The
registered NCZ neighborhood describes low activity of all channels; the ECZ
neighborhood describes non-zero activity with low resulting efficiency and a
registered destructive-interaction condition. Resulting efficiency measures the share
of channel activity surviving in the net correction; destructive interaction records
the negative part of pairwise alignment. TNZ separately denotes non-zero error input
suppressed during transport to the state. A directed transport coefficient describes
one registered input vector and does not replace spectral analysis of the full Jacobian.
A block-Jacobian probe is a matrix-free observation of such transport and is not itself
evidence of speedup or global algorithmic locality.

\subsection{Relationship Between PC-TREF and PC-CATM}
\label{sec:theory-tref-catm-link}

The two levels solve different problems and must not be collapsed into one model.
PC-CATM asks how an observed correction is structurally decomposed: which channels are
active, how they are transported, and under which combinations their aggregate effect
is large, small, or cancelled. PC-TREF asks whether that observation is sufficient for
an action: can the selected features preserve the same admissible computational
decision that would be obtained with complete exact continuation?

The dependency is
\[
  \text{PC-CATM mechanism}
  \;\longrightarrow\;
  \phi_I(x)
  \;\longrightarrow\;
  \text{PC-TREF sufficiency test}
  \;\longrightarrow\;
  a\in\mathcal A.
\]
The first arrow constructs a diagnostic representation; the second checks whether
decision-relevant distinctions are lost; the third selects an action only after the
admissibility criterion is satisfied.

This separation explains the Stage~3A/3B sequence. Layer-wise gradients and
representations show that PC/BP similarity varies with depth. Profiling and SI-MA0/1
localize cost and test measurability of mechanistic components. B1/B2 then test exact
alternative computational organizations. Only afterward is it meaningful to ask
whether pre-action state is sufficient for a continuation decision.

\subsection{QWake-PC and the Experimental QWake-FP Instantiation}
\label{sec:theory-qwake}

QWake-PC is the proper name of the compute-budget control architecture. It is not a new
kind of predictive coding and does not replace PC-TREF or PC-CATM. Architecturally,
\[
  \text{PC-TREF}\;\longrightarrow\;\text{sufficiency criterion},
  \qquad
  \text{PC-CATM}\;\longrightarrow\;\text{mechanistic features},
\]
\[
  \text{criterion + features}
  \;\longrightarrow\;
  \text{QWake-PC}
  \;\longrightarrow\;
  \text{computational action}.
\]

In general, the architecture chooses from a bounded action space, for example continue
exact computation, choose a registered bounded analytic completion, or switch to an
exact fallback path. A decision is admissible only when an estimate available
\emph{before action} indicates sufficiency of the current state, while a separate exact
post-action oracle verifies that the registered required-response tolerance was not
violated. QWake-PC therefore separates estimator and reference oracle: the latter may
use counterfactual continuation for labeling and audit but cannot be a pre-decision
feature.

The dissertation evaluates only QWake-FP, a bounded shadow instantiation for the
registered FixedPred case with $\eta=1$ and the canonical Stage~2 executor. The
registered early action belongs to the \code{ANALYTIC_COMPLETION} family and has
identifier \code{fixedpred_eta1_wavefront_completion_v1}. It replaces the remaining
canonical iterative sweeps with a bounded analytic completion and therefore still
performs computation. The complete exact suffix
\code{complete_suffix_stage2_baseline_v1} is reference and fallback. In shadow mode,
both branches execute on disposable forks of one state to produce evidence; the rule
only proposes an action and does not activate it in the live trajectory.

The PC-TREF objects map to QWake-FP as follows:
\begin{center}
\small
\begin{tabular}{p{0.27\textwidth}p{0.64\textwidth}}
\toprule
PC-TREF object & QWake-FP realization \\
\midrule
$x$ & state $x(\rho)$ contained in C1 record $\rho$ and available before action \\
$\phi_I(x)$ & allowed pre-action C2 rule features \\
$\mathcal A$ & continue the exact suffix or propose the registered analytic completion \\
$q_R^{\mathrm{QW}}$ & \texttt{EARLY\_ADMISSIBLE}/\texttt{CONTINUE\_REQUIRED}, determined by comparing the analytic candidate's required response with the post-action exact-suffix required response \\
$\delta_R$ & registered tolerance on required final response \\
$\mathfrak D_{I\to R}^{q}$ & diagnostic merging of states from different registered required classes \\
$\mathfrak D_{\pi,a_{\mathrm{early}}}$ & dangerous accept: the rule proposes the early action for a state that the oracle requires to continue \\
action cost & frozen C2 full decision-cost vector \\
\bottomrule
\end{tabular}
\end{center}
This mapping is experimental and does not enlarge PC-TREF. C2 tests a finite family of
scalar representations, not every possible $\phi_I$. In particular, it tests selective
sufficiency of the early action on accepted records, not full sufficiency
$E_I|_\Omega\subseteq E_R|_\Omega$ of the whole diagnostic representation.

C1 and C2 follow from this separation. C1 collects complete trajectories and
counterfactual post-action sufficiency labels without selecting a rule. C2 operates
only on the sealed C1 surface and tests a finite rule family using exclusively
pre-action features. The final branch therefore asks, in strict order:
\begin{enumerate}
  \item whether preterminal states exist for which the required response of the registered analytic completion is equivalent to that of the complete exact suffix within the registered tolerance;
  \item whether a non-zero subset can be recognized on the frozen calibration surface without future-result access and with zero observed dangerous accepts;
  \item whether the avoided computation remains beneficial after observation, decision, analytic-completion, and exact-fallback costs are accounted for.
\end{enumerate}

For the fixed FixedPred case, the complete canonical suffix is the exact reference for
the registered analytic candidate. Post hoc, one can define the \emph{minimal stably
admissible FixedPred prefix}: the first inference-cycle index from which the candidate's
required response remains equivalent to the exact-suffix required response for all
subsequent decision epochs. This quantity is unavailable to the rule at decision time
because it requires post-action oracle checks. A separate Rosenbaum analytic wavefront
control is used as a positive control for the special FixedPred $\eta=1$ structure,
but it does not itself authorize \code{ANALYTIC_COMPLETION} and must not be conflated
with \code{fixedpred_eta1_wavefront_completion_v1}.

QWake-FP shadow decisions use D/U/S semantics. \code{DONE} proposes the registered
analytic completion instead of the next canonical sweep; \code{UNKNOWN} means
available features are insufficient for that proposal; \code{SWEEP} continues the
canonical inference cycle. A \emph{dangerous DONE} occurs when \code{DONE} is proposed
but the candidate required response fails the registered comparison with the exact
suffix. Conversely, when the candidate is oracle-equivalent but the cheap diagnostic
returns \code{UNKNOWN}, there is a \emph{diagnostic observability gap}. QWake thus
separates inadmissibility of an action from inability to recognize existing
admissibility before action.

In the broader QWake-PC architecture, these relations can be expressed through an
operational sufficiency boundary. A reference sufficiency margin is computed after
action relative to the exact reference, while a sufficiency-boundary estimator uses
only pre-action features. A predicted sufficiency horizon is stronger: it estimates
how long the current margin may persist. C1/C2 does not test that horizon; it is bounded
to recognition of current admissibility of the registered analytic action, so the
finite scalar-rule family is not assigned a stronger capability than was tested.

The economic decision is represented by a \emph{cost vector}, not one implicit scalar.
Compute, latency, memory, diagnostic mechanism, observer, control-loop, and fallback
costs are kept distinct. Pareto admissibility can be used where multiple components
cannot be justifiedly reduced to one number. C2 used one concrete preregistered full
decision-cost accounting model, so its negative economic result is not equivalent to
impossibility of a cheap minimal controller.

A positive answer at one level does not imply a positive answer at the next. Existence
of task-admissible states does not guarantee their recognizability, and zero observed
dangerous accepts on the calibration surface establishes neither economic benefit nor
population-level dangerous-accept safety. This three-level logic is the theoretical
basis of RQ3.

\section{Adaptive Computation and Early Termination of Computation}
\label{sec:related-adaptive-computation}

Input-dependent computation predates predictive coding. Adaptive Computation Time
(ACT) lets a recurrent network adapt the number of internal steps before emitting a
result \cite{graves2016adaptive}, making compute budget part of the model's decision
rather than a fixed architectural constant.

For feed-forward networks, early termination is often implemented through intermediate
exits. Bolukbasi et al. formulate adaptive network evaluation as selecting which
components must be computed for a particular example and train an early-exit rule with
the explicit goal of reducing runtime \cite{bolukbasi2017adaptive}. Multi-exit
architectures remain active: Kubaty et al. analyze how joint or separate training of
the backbone and exit heads changes model properties \cite{kubaty2025multiexit}.

For predictive coding specifically, Zniber et al. propose EE-PCN, a dynamic early-exit
architecture for PC cycles in which a classifier evaluates quality after a cycle and
classifier confidence is compared with a user threshold; separate classifiers are used
for up to $T$ cycles \cite{zniber2025earlyexit}. This is a direct neighbor of QWake
because the controlled quantity is the number of PC cycles.

The formulations nevertheless differ. EE-PCN designs and trains an architecture with
early classifiers and a confidence threshold. QWake-FP introduces no new classifier
head per step; it studies an existing FixedPred trajectory, asks whether a bounded
pre-action state can recognize sufficiency labeled by exact post-action continuation,
and then separately accounts for full decision cost. QWake-FP is therefore not claimed
as the first PC early-exit method. Its contribution is the registered separation of
task admissibility, observed zero-danger calibration recognition, and economics.

\section{Selective Prediction, Abstention, and Risk-aware Early Exit}
\label{sec:related-selective}

Early exit creates a second question beyond saving computation: when is reduced
computation admissible with respect to output quality? In selective prediction, a
model can abstain on some inputs, creating a risk--coverage tradeoff. SelectiveNet
combines a selection function and predictive model in one trainable architecture
\cite{geifman2019selectivenet}. The important principle here is separation of coverage
and risk: a high accepted fraction is meaningless if accepted states include states
the criterion declares insufficient.

Recent early-exit work makes risk control explicit. Jazbec et al. apply risk-control
methods to early-exit networks and bind the exit decision to a predefined quality
boundary \cite{jazbec2024riskcontrol}; related work considers nested uncertainty sets
across multi-exit networks \cite{jazbec2024nested}. EERO combines early exit and
abstention with an explicit compute-budget constraint \cite{valade2025eero}.

QWake-FP is adjacent to this literature but uses a stricter bounded dangerous-accept
criterion. On calibration evidence, a rule is admissible only with zero observed
dangerous accepted decisions; non-zero coverage alone is insufficient. The economic
criterion is applied only afterward. QWake-FP can therefore be viewed as a special
case of selective computation: acceptance chooses the registered bounded analytic
completion instead of the canonical iterative suffix, while abstention continues the
exact suffix in a risk-aware fallback sense. Acceptance does not mean zero further
computation. This analogy is conceptual; QWake-FP does not inherit statistical risk
or conformal risk-control bounds that were not implemented and tested in its own protocol.

\section{Positioning of the Present Work}
\label{sec:related-positioning}

The literature provides several distinct axes that this dissertation connects in one
sequential experimental line:
\begin{enumerate}
  \item PC may approximate or, under special conditions, reproduce selected BP updates, but the applicability domain depends on the concrete algorithmic regime \cite{whittington2017approximation,song2020exact,salvatori2021convolutional,millidge2022arbitrary,rosenbaum2022relationship,rosenbaum2025correction};
  \item representation similarity requires separate diagnostics and does not automatically follow from final-error similarity \cite{raghu2017svcca,morcos2018insights,kornblith2019similarity};
  \item computational price must be measured separately from functional relation to BP \cite{zahid2023critical};
  \item adaptive computation trades completeness against cost, while selective-prediction and early-exit work adds an explicit risk--coverage dimension \cite{graves2016adaptive,bolukbasi2017adaptive,zniber2025earlyexit,geifman2019selectivenet,jazbec2024riskcontrol,valade2025eero}.
\end{enumerate}

The dissertation therefore does not offer another undifferentiated PC/BP equivalence
claim and is not reducible to a new multi-exit architecture. It sequentially tests
behavior, internal representations and gradients, mechanism, and measured cost within
one versioned experimental line, then asks the separate bounded question of whether a
registered analytic completion can replace part of an already existing iterative
computation. That separation motivates RQ1--RQ3 and the subsequent chapters.
